\documentclass[11pt,a4paper]{article}

\usepackage[T1]{fontenc}
\usepackage[utf8]{inputenc}
\usepackage{amsmath,amssymb,amsthm}
\usepackage{mathtools}
\usepackage[margin=1.1in]{geometry}
\usepackage{hyperref}
\usepackage{booktabs}

\theoremstyle{plain}
\newtheorem{theorem}{Theorem}[section]
\newtheorem{proposition}[theorem]{Proposition}
\newtheorem{lemma}[theorem]{Lemma}
\newtheorem{corollary}[theorem]{Corollary}
\theoremstyle{definition}
\newtheorem{definition}[theorem]{Definition}
\newtheorem{remark}[theorem]{Remark}
\newtheorem{question}[theorem]{Question}

\DeclareMathOperator{\Gal}{Gal}
\DeclareMathOperator{\Res}{Res}
\DeclareMathOperator{\disc}{disc}
\DeclareMathOperator{\Aut}{Aut}
\DeclareMathOperator{\Frob}{Frob}
\DeclareMathOperator{\Mon}{Mon}
\newcommand{\PP}{\mathbb{P}}
\newcommand{\FF}{\mathbb{F}}
\newcommand{\ZZ}{\mathbb{Z}}
\newcommand{\Fq}{\FF_q}
\newcommand{\Fbar}{\overline{\FF}_q}
\newcommand{\Sm}{S_m}
\newcommand{\eps}{\varepsilon}

\title{Conservation and Specialization in \\ Semaev-Based Index Calculus over Prime Fields}
\author{}
\date{}

\begin{document}
\maketitle

\begin{abstract}
Let $E$ be an elliptic curve over a field $K$ and let $S_m$ be Semaev's $m$-th
summation polynomial. We determine the Galois group of $S_m(t_1,\dots,t_{m-1},T)$
over $K(t_1,\dots,t_{m-1})$ for every $m \ge 3$: it is elementary abelian of
order $2^{m-2}$, acting simply transitively on the $2^{m-2}$ roots, and it is
the same group for every elliptic curve, every base field and every
characteristic. The proof identifies the splitting field with the function field
of $E^{m-1}/\langle -1\rangle$, where $-1$ acts diagonally; the group is the
deck group $(\ZZ/2)^{m-1}$ of $E^{m-1} \to (E/{\pm})^{m-1}$ modulo the diagonal.
Geometric and arithmetic monodromy coincide.

Two consequences are recorded. First, an exact factorization law over a finite
field $\Fq$: at every specialization outside an explicitly described degenerate
locus of size $O_m(q^{m-2})$, $S_m(a_1,\dots,a_{m-1},T)$ either splits into
$2^{m-2}$ distinct linear factors --- exactly when the $m-1$ fibres
$\pi^{-1}(a_i)$ of the $x$-map are simultaneously rational or simultaneously
irrational --- or is a product of $2^{m-3}$ distinct irreducible quadratics. No
irreducible factor of degree $\ge 3$ ever occurs, and the factor degrees are
always equal. Second, the factor-base locus, on which index-calculus relation
search operates, lies entirely inside the totally split locus, at every $m$.

Combined with an elementary conservation identity --- the total number of
$m$-term decompositions over a factor base of fixed size is a binomial
coefficient, so mean decomposition yield is independent of factor-base geometry
--- these give a coherent negative statement: generic-fibre arithmetic statistics
and factor-base relation statistics are different objects, and several proposed
relation-rate levers are fixed either by conservation or by complete splitting on
the rational-point locus. Nothing here claims an ECDLP speedup, a security
reduction, or a new hardness result; several of the ingredients are classical and
are marked as such. All statements are verified computationally on $20$ curves
over $6$ primes at $m = 3, 4, 5$, including $257\,061$ exhaustive specializations
with zero discrepancies, against null controls that the same instrument rejects.
\end{abstract}

\tableofcontents

%======================================================================
\section{Introduction}

Semaev's summation polynomials \cite{semaev2004} translate the statement
``$m$ points of an elliptic curve sum to the identity'' into a polynomial
condition on their $x$-coordinates alone, and they are the algebraic engine of
index-calculus attacks on the elliptic curve discrete logarithm problem
(ECDLP) \cite{gaudry2009,diem2011,fghr2014}. A recurring hope in that line of
work is that some \emph{choice} --- of curve, of curve family, of factor base ---
might raise the rate at which relations are found. This paper makes three
statements about which choices can possibly matter, and finds that a specific,
identifiable class of them cannot.

The main object is the cover of parameter space cut out by $S_m$ in its last
variable. Write $n = m-1$, let $t_1,\dots,t_n$ be independent transcendentals
over $K$, and consider
\[
  S_m(t_1,\dots,t_n,T) \in K(t_1,\dots,t_n)[T],
\]
a polynomial of degree $2^{m-2}$ in $T$. Theorem~\ref{thm:main} says its Galois
group is $(\ZZ/2)^{m-2}$ in its regular representation, for every $E$, every $K$
and every characteristic, with geometric monodromy equal to arithmetic monodromy.

The proof is short once the right object is named. The roots of
$S_m(t_1,\dots,t_n,T)$ are the $x$-coordinates of the signed sums
$\eps_1 P_1 + \cdots + \eps_n P_n$ of the tautological points $P_i$ with
$x(P_i) = t_i$ --- a classical observation, immediate from the defining property
of $S_m$. The splitting field is therefore contained in $K(E^n)$, which is Galois
over $K((E/{\pm})^n) = K(t_1,\dots,t_n)$ with group $(\ZZ/2)^n$, the deck group of
$E^n \to (E/{\pm})^n$; no computation with square classes is needed, and the
argument is insensitive to the characteristic. The diagonal $-1$ fixes every root
because $x(R) = x(-R)$, and --- this is the only step with content --- nothing else
does, because a further coincidence among the roots would force a nontrivial
relation $\sum_{i \in S} \eps_i P_i \in E[2]$ at the generic point of $E^n$, which
is a proper closed condition. So the group is $(\ZZ/2)^n$ modulo the diagonal.
Because the group is abelian, the stabilizer of a root is normal, so the root
field is already the splitting field: the cover is Galois.

Over a finite field this yields an exact factorization law
(Theorem~\ref{thm:frob}). Frobenius acts on the deck group through the
$(m-1)$-tuple recording, for each $i$, whether the fibre $\pi^{-1}(a_i)$ of the
$x$-map is $\Fq$-rational. Because the action of the group on the roots is
regular, every nonidentity element is fixed-point-free, so the fibre is either
totally split or has no rational point at all: the dichotomy in the abstract.
This is a sharp, immediately checkable statement about a heavily studied family
of polynomials.

Specializing to the locus where each $a_i$ is the $x$-coordinate of a
\emph{rational} point --- the factor-base locus --- all the rationality invariants
agree and the fibre splits completely, at every $m$
(Theorem~\ref{thm:factorbase}). This is the algebraic-geometry restatement of a
fact index calculus already uses operationally, and we are careful in
\S\ref{sec:notalgo} to explain why it is not an algorithm. Its interest is
negative and methodological: a Frobenius/Chebotarev census of the generic fibre
measures a statistic that is \emph{identically constant} on the locus where
relation search actually operates, so such a census cannot detect a
relation-rate anomaly there --- there is none to detect.

Section~\ref{sec:conservation} records the complementary elementary fact on the
combinatorial side: a double count fixes the mean number of $m$-term
decompositions per target at $\binom{B+m-1}{m}/N$ for \emph{every} factor base of
size $B$ in a group of order $N$. Mean yield is not a design lever; only the
distribution is free.

\subsection*{What this paper does not claim}

It claims no improvement to any ECDLP algorithm, no security reduction, no new
hardness assumption, and no change to any complexity estimate. Every asymptotic
statement about index calculus in the literature stands untouched. The results
are structural facts about a family of polynomials, plus one combinatorial
identity, and their content for cryptanalysis is exclusively that of ruling
certain proposals out.

It also claims no novelty for its classical ingredients. The identification of
the roots as signed sums, the degree $2^{m-2}$, the symmetry, the absolute
irreducibility of $S_m$, the resultant recursion, and the operational
``solve for the last variable'' step are all standard
\cite{semaev2004,gaudry2009,ky2015}. What is offered as possibly new is the
Galois-theoretic identification for all $m$ (Theorem~\ref{thm:main}), the exact
factorization law (Theorem~\ref{thm:frob}), and the synthesis in
\S\ref{sec:synthesis}. \S\ref{sec:novelty} states honestly what search was and
was not possible from the environment in which this work was done, and does not
assert novelty beyond what that search supports.

%======================================================================
\section{Notation and the summation cover}\label{sec:setup}

Throughout, $K$ is a field and $E/K$ is an elliptic curve with origin $O$ and
inversion $\iota = [-1]$. Let
\[
  \pi \colon E \longrightarrow E/\langle\iota\rangle \cong \PP^1_K
\]
be the quotient map. We use only the following two properties of $\pi$, both of
which hold in every characteristic:

\begin{itemize}
\item[(P1)] $\pi$ is finite of degree $2$ and separable, and
  $K(E)/K(\PP^1)$ is Galois with group $\langle \iota \rangle \cong \ZZ/2$.
\item[(P2)] $\iota \neq \mathrm{id}_E$.
\end{itemize}

(P2) is immediate since $E$ has points of order $>2$. For (P1): if
$\operatorname{char} K \neq 2$ we may take $E : y^2 = f(x) = x^3 + Ax + B$ with
$f$ separable and $\pi = x$, and $\iota(x,y) = (x,-y)$; if
$\operatorname{char} K = 2$ then $\iota(x,y) = (x, y + a_1 x + a_3)$ with
$(a_1,a_3) \neq (0,0)$, which is again a separable involution fixing $x$. In both
cases $[K(E):K(x)] = 2$ and $\iota$ generates the automorphism group of the
extension.

\begin{definition}[Summation polynomials]\label{def:sm}
$S_2(X_1,X_2) = X_1 - X_2$, and for $m \ge 3$, $S_m \in K[X_1,\dots,X_m]$ is the
reduced defining polynomial of the Zariski closure of
\[
  \bigl\{ (\pi(Q_1),\dots,\pi(Q_m)) \;:\; Q_i \in E(\overline{K}),\;
      Q_1 + \cdots + Q_m = O \bigr\} \subseteq (\PP^1)^m .
\]
Equivalently (\cite{semaev2004}; see also \cite[Prop.~2.1]{ky2015}), for
$\xi_1,\dots,\xi_m \in \overline{K}$ one has $S_m(\xi_1,\dots,\xi_m) = 0$ if and
only if there are $Q_i \in E(\overline{K})$ with $\pi(Q_i) = \xi_i$ and
$\sum_i Q_i = O$. For $\operatorname{char} K \neq 2,3$ and
$E : y^2 = x^3 + Ax + B$,
\begin{align}
  S_3(X_1,X_2,X_3) &= (X_1-X_2)^2 X_3^2
     - 2\bigl((X_1+X_2)(X_1X_2+A) + 2B\bigr) X_3 \notag \\
   &\qquad + \bigl( (X_1X_2 - A)^2 - 4B(X_1+X_2) \bigr),
   \label{eq:s3}
\end{align}
and $S_m$ for $m \ge 4$ is obtained from the resultant recursion
\begin{equation}\label{eq:recursion}
  S_m(X_1,\dots,X_m) =
  \Res_X\bigl( S_{m-1}(X_1,\dots,X_{m-2},X),\; S_3(X_{m-1},X_m,X) \bigr).
\end{equation}
$S_m$ is symmetric and absolutely irreducible, of degree $2^{m-2}$ in each
variable for $m \ge 3$ \cite{semaev2004,gaudry2009}.
\end{definition}

We now name the cover. Fix $m \ge 3$ and put $n := m-1 \ge 2$.

\begin{definition}[The summation cover]\label{def:cover}
Let $t_1,\dots,t_n$ be independent transcendentals over $K$ and set
\[
  L := K(t_1,\dots,t_n) = K\bigl((\PP^1)^n\bigr), \qquad
  M := K(E^n).
\]
Let $\Gamma := \langle \iota \rangle^n \cong (\ZZ/2)^n$ act on $E^n$
coordinatewise, let $\Delta := \langle (\iota,\dots,\iota) \rangle \cong \ZZ/2$
be its diagonal subgroup, and write $\sigma_i \in \Gamma$ for the generator
acting by $\iota$ in the $i$-th coordinate only.
Let $s \colon E^n \to E$, $(Q_1,\dots,Q_n) \mapsto Q_1 + \cdots + Q_n$, be the
summation morphism and put $\varphi := \pi \circ s \colon E^n \to \PP^1$.
Finally set
\[
  X_n := E^n / \Delta, \qquad \rho \colon X_n \to (\PP^1)^n
\]
the map induced by $\pi^n$ (which is $\Gamma$-invariant, hence $\Delta$-invariant).
\end{definition}

We identify $L$ with $K(t_1,\dots,t_n)$ via $t_i = \pi \circ \mathrm{pr}_i$, and
write $P_i \in E(M)$ for the $i$-th tautological point, i.e. the $i$-th
projection of the generic point of $E^n$; thus $\pi(P_i) = t_i$ and
$\{ \pm P_i \}$ is exactly the fibre of $\pi$ over $t_i$.

\begin{lemma}\label{lem:gamma}
$M/L$ is Galois with $\Gal(M/L) = \Gamma \cong (\ZZ/2)^n$.
\end{lemma}

\begin{proof}
$\Gamma$ acts faithfully on $E^n$ by (P2): if $(\eps_1,\dots,\eps_n)$ acts
trivially then each $\eps_i$ acts trivially on $E$, so $\eps_i = \mathrm{id}$.
Hence by Artin's theorem $M/M^{\Gamma}$ is Galois with group $\Gamma$. Finally
$M^{\Gamma} = K(E^n/\Gamma) = K\bigl((E/\langle\iota\rangle)^n\bigr) = K((\PP^1)^n) = L$,
using $E^n/\Gamma = (E/\langle \iota\rangle)^n$ and (P1).
\end{proof}

Note what Lemma~\ref{lem:gamma} replaces. In coordinates one would argue that
$f(t_1),\dots,f(t_n)$ are independent in $L^{\times}/(L^{\times})^2$; that
argument is correct for $\operatorname{char} K \neq 2$ but needs the Weierstrass
model and a squarefreeness check. The fibre-product formulation gives the same
conclusion with no computation and in every characteristic.

For $\eps \in \{\pm 1\}^n$ write
\begin{equation}\label{eq:root}
  \theta_{\eps} \;:=\; \pi\bigl( \eps_1 P_1 + \cdots + \eps_n P_n \bigr)
  \;\in\; \PP^1(M),
\end{equation}
so that $\theta_{-\eps} = \theta_{\eps}$, and $\sigma \in \Gamma$ acts by
$\sigma(\theta_{\eps}) = \theta_{\sigma \eps}$ where $\Gamma = \{\pm1\}^n$ acts on
$\{\pm1\}^n$ by coordinatewise multiplication.

%======================================================================
\section{The structure theorem}\label{sec:main}

\subsection{The two lemmas}

\begin{lemma}[Root identification]\label{lem:roots}
Let $\Omega \supseteq M$ be an algebraically closed field and $\theta \in \Omega$.
Then $S_m(t_1,\dots,t_n,\theta) = 0$ if and only if $\theta = \theta_{\eps}$ for
some $\eps \in \{\pm1\}^n$. Moreover $\theta_{\eps} = \infty$ if and only if
$\eps_1 P_1 + \cdots + \eps_n P_n = O$.
\end{lemma}

\begin{proof}
By Definition~\ref{def:sm} applied over $\Omega$, $S_m(t_1,\dots,t_n,\theta) = 0$
iff there exist $Q_1,\dots,Q_n, R \in E(\Omega)$ with $\pi(Q_i) = t_i$,
$\pi(R) = \theta$ and $Q_1 + \cdots + Q_n + R = O$. The fibre of $\pi$ over $t_i$
is $\{\pm P_i\}$, so $Q_i = \eps_i P_i$ for some sign, and then
$R = -\sum_i \eps_i P_i$, whence $\theta = \pi(R) = \theta_{\eps}$. The converse
is the same computation read backwards. The last sentence is
$\pi^{-1}(\infty) = \{O\}$.
\end{proof}

\begin{lemma}[Generic separability]\label{lem:sep}
The $2^{n-1}$ values $\theta_{\eps}$, $\eps \in \{\pm1\}^n/\{\pm1\}$, are pairwise
distinct, and none of them equals $\infty$.
\end{lemma}

\begin{proof}
For $\emptyset \neq S \subseteq \{1,\dots,n\}$ and $\eps \in \{\pm 1\}^{S}$ let
\[
  a_{S,\eps} \colon E^n \longrightarrow E, \qquad
  (Q_1,\dots,Q_n) \longmapsto \sum_{i \in S} \eps_i Q_i .
\]
Each $a_{S,\eps}$ is a surjective homomorphism of group varieties: already its
restriction to the $i$-th factor for one $i \in S$, namely $Q \mapsto \eps_i Q$,
is surjective. Hence for any closed subscheme $Z \subsetneq E$ of dimension $0$
the preimage $a_{S,\eps}^{-1}(Z)$ is closed of dimension $n-1$, so it is a proper
closed subscheme of $E^n$.

Now suppose $\theta_{\eps} = \theta_{\eps'}$ with $\eps' \not\equiv \pm\eps$. By
\eqref{eq:root} this means $\sum_i \eps_i P_i = \pm \sum_i \eps'_i P_i$. In the
$+$ case, subtracting gives $\sum_{i \in S} 2\eps_i P_i = O$ with
$S = \{ i : \eps_i \neq \eps'_i \}$, i.e. $\sum_{i \in S} \eps_i P_i \in E[2]$; in
the $-$ case, adding gives the same conclusion with $S$ replaced by its
complement. Since $\eps' \neq \eps$ and $\eps' \neq -\eps$, both $S$ and its
complement are nonempty. Likewise $\theta_{\eps} = \infty$ means
$a_{\{1,\dots,n\},\eps}$ sends the generic point to $O$.

So each failure would place the generic point of $E^n$ in one of the finitely
many proper closed subschemes $a_{S,\eps}^{-1}(E[2])$ (for $\emptyset \neq S
\subsetneq \{1,\dots,n\}$) or $a_{\{1,\dots,n\},\eps}^{-1}(O)$. As $E^n$ is
integral, its generic point lies in no proper closed subscheme.
\end{proof}

\begin{remark}
$E[2]$ is a finite group scheme of order $4$ in every characteristic --- in
characteristic $2$ it is non-reduced --- and Lemma~\ref{lem:sep} only uses that it
is finite, so no characteristic is excluded.
\end{remark}

\subsection{The theorem}

\begin{theorem}\label{thm:main}
Let $K$ be any field, $E/K$ an elliptic curve, $m \ge 3$ and $n = m-1$. Then, with
the notation of \S\ref{sec:setup}:
\begin{enumerate}
\item $S_m(t_1,\dots,t_n,T)$ is separable of degree $2^{m-2}$ over
  $L = K(t_1,\dots,t_n)$, with root set $\{ \theta_{\eps} \}$.
\item Its splitting field over $L$ is $N = M^{\Delta} = K(X_n)$, the function
  field of $E^{n}/\Delta$.
\item $N/L$ is Galois with
  \[
    \Gal\bigl( S_m(t_1,\dots,t_n,T) / L \bigr) \;=\; \Gamma/\Delta
      \;\cong\; (\ZZ/2)^{m-2},
  \]
  acting simply transitively on the $2^{m-2}$ roots. In particular the cover
  $\rho \colon X_n \to (\PP^1)^n$ is Galois.
\item $K$ is algebraically closed in $N$; consequently the arithmetic and
  geometric monodromy groups agree,
  \[
    \Mon^{\mathrm{arith}} \;=\; \Mon^{\mathrm{geom}} \;\cong\; (\ZZ/2)^{m-2}.
  \]
\end{enumerate}
All four statements are independent of $E$, of $K$ and of $\operatorname{char} K$.
\end{theorem}

\begin{proof}
(1) By Lemma~\ref{lem:roots} the roots of $S_m(t_1,\dots,t_n,T)$ in an algebraic
closure of $M$ are among the $\theta_{\eps}$, of which there are at most
$2^{n-1} = 2^{m-2}$ distinct ones; by Lemma~\ref{lem:sep} there are exactly that
many, all finite. Since $\deg_T S_m = 2^{m-2}$ (Definition~\ref{def:sm}), these
account for all roots with multiplicity one, so $S_m(t_1,\dots,t_n,T)$ is
separable and its roots are exactly the $\theta_{\eps}$.

(2), (3) By Lemma~\ref{lem:gamma}, $\Gal(M/L) = \Gamma$, and by \eqref{eq:root}
$\Gamma$ permutes the roots by $\sigma \cdot \theta_{\eps} = \theta_{\sigma\eps}$.
The kernel of this permutation action is
\[
  \{ \sigma \in \{\pm1\}^n : \theta_{\sigma\eps} = \theta_{\eps}
      \text{ for all } \eps \}
  = \{ \sigma : \sigma\eps \equiv \pm \eps \text{ for all } \eps \}
  = \Delta,
\]
where the middle equality uses Lemma~\ref{lem:sep} (distinctness) and the last one
is immediate: $\sigma\eps = \eps$ for all $\eps$ forces $\sigma = (1,\dots,1)$, and
$\sigma\eps = -\eps$ for all $\eps$ forces $\sigma = (-1,\dots,-1)$. Hence
$\Gamma$ acts on the $2^{m-2}$ roots through the faithful quotient $\Gamma/\Delta$,
which has order $2^{m-2}$ and is therefore \emph{regular} (simply transitive) on
them.

Let $\theta := \theta_{(1,\dots,1)}$. Its stabilizer in $\Gamma$ is $\Delta$, so
$[L(\theta):L] = [\Gamma : \Delta] = 2^{m-2}$ and $L(\theta) = M^{\Delta}$. Since
$\Gamma$ is abelian, $\Delta \trianglelefteq \Gamma$, so $M^{\Delta}$ is Galois
over $L$; therefore $M^{\Delta}$ already contains all conjugates of $\theta$, i.e.
all roots, and it is generated by $\theta$. Thus $N = L(\theta) = M^{\Delta} =
K(E^n/\Delta) = K(X_n)$ and $\Gal(N/L) = \Gamma/\Delta \cong (\ZZ/2)^{m-2}$.

(4) $E$ is geometrically integral, hence so is $E^n$, hence so is its quotient
$X_n = E^n/\Delta$ by the finite group $\Delta$. A geometrically integral
$K$-variety has $K$ algebraically closed in its function field, so
$K \cap N = K$ and the constant field extension of $N/L$ is trivial. Therefore
$\Mon^{\mathrm{geom}} = \Gal(N\overline{K}/L\overline{K}) = \Gal(N/L) =
\Mon^{\mathrm{arith}}$.
\end{proof}

\begin{remark}[Why the group cannot depend on the curve]\label{rem:whyfixed}
Statement (4) is the operative one for cryptanalysis. Arithmetic monodromy
strictly larger than geometric monodromy is exactly the situation in which the
splitting behaviour of a cover can depend on the base field --- and hence, over a
finite field, on the curve. Theorem~\ref{thm:main}(4) says that situation does
not arise here: the entire cover is defined by the deck group of
$E^n \to (E/{\pm})^n$, which is $(\ZZ/2)^n$ for structural reasons that see
nothing about $E$ beyond the fact that $\iota \neq \mathrm{id}$.
\end{remark}

\subsection{Corollaries}

\begin{corollary}[Irreducibility]\label{cor:irred}
$S_m(t_1,\dots,t_n,T)$ is irreducible over $L$ and over $\overline{K}(t_1,\dots,t_n)$.
\end{corollary}

\begin{proof}
Transitivity of the Galois action, plus Theorem~\ref{thm:main}(4) for the
geometric statement.
\end{proof}

This recovers, rather than improves on, the known absolute irreducibility of
$S_m$ \cite{semaev2004,gaudry2009}; it is recorded because it is the same
statement read off the group.

\begin{corollary}[Imprimitivity and block systems]\label{cor:blocks}
For $m = 3$ the group is $\ZZ/2$ acting on $2$ points, which is primitive for
degree reasons. For $m \ge 4$ the group is abelian and regular of composite
degree $2^{m-2} \ge 4$, hence \emph{imprimitive}; its systems of blocks are in
bijection with the subgroups $H \le (\ZZ/2)^{m-2}$, the blocks being the cosets
of $H$.

For each $\emptyset \neq I \subsetneq \{1,\dots,n\}$ the map
$\theta_{\eps} \mapsto \pi\bigl(\sum_{i \in I} \eps_i P_i\bigr)$ is well defined on
roots and its level sets form the block system attached to
$H_I = \{ \sigma \in \{\pm1\}^n : \sigma|_I = \pm \mathbf{1} \}/\Delta$. These are
the block systems exhibited by the resultant recursion \eqref{eq:recursion}. For
$m \ge 5$ they are not all of them: $(\ZZ/2)^{m-2}$ has subgroups not of the form
$H_I$.
\end{corollary}

The point of Corollary~\ref{cor:blocks} is deflationary. Imprimitivity of a
solving system is sometimes a lever --- it can permit a resolvent of smaller degree
--- but here the block systems are exactly the sign-subgroup systems already
visible in the recursion that defines $S_m$, and the associated resolvents are
summation polynomials of smaller arity. Nothing new is exposed.

\begin{corollary}[Discriminant square class]\label{cor:disc}
For $m = 3$, $\disc_T S_3(t_1,t_2,T) = 16\,f(t_1)f(t_2)$ (in
$\operatorname{char} \neq 2,3$), which is not a square in $L$. For every
$m \ge 4$, $\disc_T S_m(t_1,\dots,t_n,T)$ \emph{is} a square in $L$.
\end{corollary}

\begin{proof}
The $m=3$ identity is an identity in $\ZZ[X_1,X_2,A,B]$; it follows from
\eqref{eq:s3} by expansion, or conceptually: the two roots are
$x(P_1 + P_2)$ and $x(P_1 - P_2)$, and the chord formula gives
$x(P_1+P_2) - x(P_1-P_2) = -4y_1y_2/(t_1-t_2)^2$, whence
$\disc = c_2^2 (\theta_+ - \theta_-)^2 = 16 y_1^2 y_2^2 = 16 f(t_1)f(t_2)$ with
$c_2 = (t_1-t_2)^2$ the leading coefficient. It is not a square because
$f(t_1)f(t_2)$ is squarefree in the unique factorization domain
$\overline{K}[t_1,t_2]$ (its six irreducible factors are pairwise distinct
precisely because $E$ is nonsingular).

For $m \ge 4$: the discriminant of a separable polynomial is a square in the base
field if and only if the Galois group is contained in the alternating group. By
Theorem~\ref{thm:main} the group acts regularly, so every nonidentity element is a
fixed-point-free involution on $2^{m-2}$ points, i.e. a product of $2^{m-3}$
transpositions. For $m \ge 4$ we have $2^{m-3} \ge 2$, an even number, so every
element is even. For $m = 3$, $2^{m-3} = 1$ is odd, which is the exception.
\end{proof}

Corollary~\ref{cor:disc} is worth stating because it locates the $m=3$
discriminant identity precisely: it is not the first case of a family but the
unique case in which the discriminant carries information. For $m \ge 4$ the
quadratic character of $\disc_T S_m$ is constant, and any statistic built on it is
vacuous.

\subsection{The branch locus}\label{sec:branch}

Lemma~\ref{lem:sep} also describes exactly where the cover degenerates. Over a
point $(a_1,\dots,a_n)$ of the base, the fibre fails to consist of $2^{m-2}$
distinct affine points precisely when one of the closed conditions in the proof
holds. Explicitly:

\begin{proposition}\label{prop:branch}
Let $a = (a_1,\dots,a_n) \in \overline{K}^n$ and $P_i$ a point with
$\pi(P_i) = a_i$. Then
\begin{enumerate}
\item $\deg_T S_m(a_1,\dots,a_n,T) < 2^{m-2}$ (a root at infinity) if and only if
  $S_{n}(a_1,\dots,a_n) = 0$, i.e. some signed sum $\sum_i \eps_i P_i$ equals $O$;
\item two roots collide if and only if $\sum_{i \in S} \eps_i P_i \in E[2]$ for
  some $\emptyset \neq S \subsetneq \{1,\dots,n\}$ and some signs. The case
  $S = \{i\}$ is $P_i \in E[2]$, i.e. $a_i$ is a branch point of $\pi$.
\end{enumerate}
The union of these loci is a proper closed subset of $(\PP^1)^n$ of codimension
$1$, so it has $O_m(q^{n-1})$ points over $\Fq$.
\end{proposition}

For $m = 3$ this is completely explicit: $n = 2$, the only proper nonempty
subsets are $\{1\}$ and $\{2\}$, so the degenerate locus is
$\{t_1 = t_2\} \cup \{ f(t_1)f(t_2) = 0 \}$, the first component being
$S_2(t_1,t_2) = t_1 - t_2 = 0$. Over $\FF_p$ its point count is
$p + Z(2p - Z) - Z$ where $Z = \#\{ x : f(x) = 0\} \in \{0,1,3\}$, which is
exactly what is measured in \S\ref{sec:verify}.

\subsection{Special $j$-invariants, complex multiplication, and characteristic}
\label{sec:special}

Three families are worth checking explicitly, because they are the standard
sources of exceptional behaviour and because a claim of the form ``for every
curve'' must survive them.

\paragraph{Extra automorphisms ($j = 0$, $j = 1728$).}
If $\operatorname{char} K \neq 2,3$ then $\Aut(E,O) = \mu_2$ except for
$j = 1728$ ($\mu_4$) and $j = 0$ ($\mu_6$). Extra automorphisms do not enter
Theorem~\ref{thm:main} at any point: the construction uses only $\iota$ and the
quotient $\pi$, and Lemma~\ref{lem:gamma} needs only (P1) and (P2). What an extra
automorphism $\zeta \in \Aut(E,O)$ does is act on $E/\langle\iota\rangle = \PP^1$,
hence on the \emph{base} $(\PP^1)^n$, commuting with $\rho$. It is a symmetry of
the cover over an automorphism of the base, not a deck transformation of $\rho$,
so it does not enlarge or shrink the monodromy group. This is verified
numerically in \S\ref{sec:verify} on $j=0$ curves with $p \equiv 1$ and
$p \equiv 2 \pmod 3$ and on $j = 1728$ curves with $p \equiv 1$ and
$p \equiv 3 \pmod 4$ --- i.e. with the extra automorphisms rational and
irrational.

\paragraph{Complex multiplication and supersingularity.}
Likewise irrelevant. $\mathrm{End}(E)$ never appears; the deck group of
$E^n \to (E/{\pm})^n$ is $(\ZZ/2)^n$ whether or not $E$ has CM, and whether or
not $E$ is supersingular. A supersingular curve is included in the verification
battery.

\paragraph{Characteristic.}
Theorem~\ref{thm:main} is stated and proved in every characteristic, because it
uses only (P1) and (P2). Three caveats are nevertheless real and are stated
rather than hidden.
\begin{itemize}
\item The \emph{Weierstrass computations} --- \eqref{eq:s3}, the explicit
  $\disc_T S_3 = 16 f(t_1) f(t_2)$, and the description of the branch locus by
  $f(t_i) = 0$ --- require $\operatorname{char} K \neq 2,3$. In characteristic $2$
  and $3$ the corresponding summation polynomials must be taken in the
  appropriate model \cite{gg2014}; the group-theoretic conclusion is unchanged
  but the formulas are not.
\item In characteristic $2$, $E[2]$ has order $2$ (ordinary) or is infinitesimal
  (supersingular) as a group scheme; Lemma~\ref{lem:sep} is unaffected, but
  Corollary~\ref{cor:disc} is not meaningful, since ``square class of the
  discriminant'' is the wrong invariant in characteristic $2$.
\item Theorem~\ref{thm:frob} below is stated over $\Fq$ using a splitting
  invariant $c(\cdot) \in \ZZ/2$ that is characteristic-free; only its
  identification with a quadratic character requires $q$ odd.
\end{itemize}

%======================================================================
\section{An exact factorization law over a finite field}\label{sec:frob}

Let $K = \Fq$ now. For $a \in \PP^1(\Fq)$ define the \emph{splitting invariant}
\begin{equation}\label{eq:c}
  c(a) := \begin{cases}
    0 & \text{if } \pi^{-1}(a) \subseteq E(\Fq) \text{ and } \#\pi^{-1}(a) = 2, \\
    1 & \text{if } \pi^{-1}(a) \cap E(\Fq) = \emptyset, \\
    \ast & \text{if } \#\pi^{-1}(a) = 1 \quad (a \text{ a branch point}).
  \end{cases}
\end{equation}
For $q$ odd and $E : y^2 = f(x)$ one has $c(a) = 0$ if $\chi(f(a)) = +1$,
$c(a) = 1$ if $\chi(f(a)) = -1$ and $c(a) = \ast$ if $f(a) = 0$, where $\chi$ is
the quadratic character of $\Fq$. Thus $c(a)$ records whether Frobenius fixes or
swaps the two points above $a$.

\begin{definition}\label{def:good}
Call $a = (a_1,\dots,a_n) \in \Fq^n$ \emph{good} if $S_n(a_1,\dots,a_n) \neq 0$
and $\sum_{i \in S} \eps_i P_i \notin E[2]$ for every
$\emptyset \neq S \subsetneq \{1,\dots,n\}$ and every choice of signs, where
$P_i \in E(\FF_{q^2})$ is either point above $a_i$. By
Proposition~\ref{prop:branch}, $a$ is good if and only if
$S_m(a_1,\dots,a_n,T)$ is separable of degree $2^{m-2}$; the non-good locus has
$O_m(q^{n-1})$ points. Note $c(a_i) \neq \ast$ for good $a$.
\end{definition}

\begin{theorem}[Factorization law]\label{thm:frob}
Let $a \in \Fq^n$ be good. Then $S_m(a_1,\dots,a_n,T) \in \Fq[T]$ is squarefree
of degree $2^{m-2}$ and exactly one of the following holds.
\begin{enumerate}
\item If $c(a_1) = c(a_2) = \cdots = c(a_n)$, then $S_m(a_1,\dots,a_n,T)$ splits
  into $2^{m-2}$ distinct linear factors over $\Fq$, with roots
  $\{ x(\eps_1 P_1 + \cdots + \eps_n P_n) \}$.
\item Otherwise $S_m(a_1,\dots,a_n,T)$ is a product of $2^{m-3}$ distinct
  irreducible quadratics over $\Fq$, and has \emph{no} root in $\Fq$.
\end{enumerate}
In particular: no irreducible factor of degree $\ge 3$ ever occurs; all
irreducible factors always have the same degree; and the factorization type
depends on $a$ only through the class of $(c(a_1),\dots,c(a_n))$ in
$(\ZZ/2)^n / \Delta$.
\end{theorem}

\begin{proof}
Goodness says the specialization lies in the étale locus of the cover
$\rho \colon X_n \to (\PP^1)^n$, so the factorization type of the fibre equals the
cycle type of the Frobenius element in $\Gal(N/L) = \Gamma/\Delta$ acting on the
$2^{m-2}$ roots (Dedekind's theorem for specializations of a Galois cover; see
e.g. \cite[Ch.~6]{fj2008}).

Frobenius acts on $M = L(E^n)$ through its action on each factor, and on the
$i$-th factor it fixes or swaps the two points of $\pi^{-1}(a_i)$ according to
$c(a_i)$: that is, the Frobenius element of $\Gamma = \{\pm1\}^n$ is
$\bigl( (-1)^{c(a_1)},\dots,(-1)^{c(a_n)} \bigr)$. Its image in $\Gamma/\Delta$ is
trivial exactly when all $c(a_i)$ agree, which is case (1); then Frobenius fixes
every root and the fibre is totally split, and the roots are the signed sums by
Lemma~\ref{lem:roots}.

Otherwise the Frobenius element is a nontrivial element of $\Gamma/\Delta$, which
acts \emph{regularly} on the roots by Theorem~\ref{thm:main}(3). A nontrivial
element of a regular permutation group has no fixed point, and here it has order
$2$, so its cycle type is $2^{2^{m-3}}$: the fibre is a disjoint union of
$2^{m-3}$ Frobenius orbits of size $2$, giving $2^{m-3}$ distinct irreducible
quadratic factors and no rational root.
\end{proof}

\begin{corollary}[Exact counts and density]\label{cor:density}
Let $S = \#\{a \in \Fq : c(a) = 0\}$, $N = \#\{a : c(a) = 1\}$ and
$Z = \#\{a : c(a) = \ast\}$, so $S + N + Z = q$ and $\#E(\Fq) = 1 + Z + 2S$.
For $\eps \in \{\pm1\}^n$ with $k$ entries equal to $+1$, the number of
$a \in \Fq^n$ whose splitting-invariant vector lies in the class of $\eps$ is
exactly
\begin{equation}\label{eq:classcount}
  S^k N^{n-k} + N^k S^{n-k}.
\end{equation}
Consequently the number of \emph{good} $a$ giving a totally split fibre is
$S^n + N^n$ minus the number of non-good $a$ with all $c(a_i)$ equal, and the
density of totally split fibres satisfies
\[
  \Bigl| \frac{S^n + N^n}{q^n} - 2^{-(m-2)} \Bigr|
    \;\le\; \frac{31 (m-1)^2}{2^{m-1} \, q}
  \qquad \text{for } q \ge 9(m-1)^2 .
\]
\end{corollary}

\begin{proof}
\eqref{eq:classcount} is immediate. For the density write $S/q = (1-\alpha)/2$ and
$N/q = (1+\beta)/2$ with $\alpha = (t+Z)/q$, $\beta = (t-Z)/q$ and $t = q+1-\#E(\Fq)$
the trace, so $|\alpha|, |\beta| \le (2\sqrt q + 3)/q =: \delta$ by Hasse. Then
\[
  \frac{S^n+N^n}{q^n} = 2^{-n}\bigl[(1-\alpha)^n + (1+\beta)^n\bigr]
   = 2^{-n}\Bigl[ 2 + n(\beta - \alpha) + \textstyle\sum_{k \ge 2}\binom{n}{k}
       \bigl((-\alpha)^k + \beta^k\bigr) \Bigr].
\]
The linear term is $n(\beta - \alpha) = -2nZ/q$, which is \emph{independent of the
trace}: the Hasse-interval contribution cancels at first order. The tail is
bounded by $2\bigl[(1+\delta)^n - 1 - n\delta\bigr] \le (n\delta)^2 e^{n\delta}$,
and $Z \le 3$, giving
$\bigl| (S^n+N^n)/q^n - 2^{1-n} \bigr| \le 2^{-n}\bigl[ 6n/q + 9n^2 e^{3n/\sqrt q}/q \bigr]$,
which for $q \ge 9n^2$ is at most $2^{-n} \cdot 31 n^2 / q$.
\end{proof}

\begin{remark}
The cancellation of the trace term is the general form of the $m = 3$ observation
that the split frequency is $\tfrac12 + O(1/q)$ rather than
$\tfrac12 + O(1/\sqrt q)$: the two error terms $\pm t/q$ enter $S^n$ and $N^n$ with
opposite signs and cancel to first order at every $n$. So the split density
carries no first-order information about the curve --- a further sense in which
this statistic is not a lever.
\end{remark}

%======================================================================
\section{Specialization to the factor-base locus}\label{sec:factorbase}

Index-calculus relation search over $\FF_p$ works with a factor base
$\mathcal{F} \subseteq E(\FF_p)$ and looks for relations
$P_1 + \cdots + P_{m-1} = R$ with $P_i \in \mathcal{F}$. The coordinates that
enter the summation polynomial are therefore $x$-coordinates of \emph{rational}
points --- a proper subvariety's worth of the parameter space, of density about
$1/2$ in each coordinate.

\begin{theorem}[Complete splitting on the factor-base locus]\label{thm:factorbase}
Let $a_1,\dots,a_n \in \Fq$ be such that each $a_i = x(P_i)$ for some
$P_i \in E(\Fq)$, and suppose $a$ is good. Then $S_m(a_1,\dots,a_n,T)$ splits into
$2^{m-2}$ distinct linear factors over $\Fq$, with root set exactly
\[
  \bigl\{\, x(\eps_1 P_1 + \cdots + \eps_n P_n) \;:\; \eps \in \{\pm1\}^n \,\bigr\}
  \;\subseteq\; \Fq,
\]
at every $m \ge 3$. Equivalently, the decomposition group of the cover
$\rho$ at every such point is trivial: the factor-base locus is contained in the
totally split locus.
\end{theorem}

\begin{proof}
$c(a_i) = 0$ for all $i$, so Theorem~\ref{thm:frob}(1) applies. Directly: the
$P_i$ are $\Fq$-rational, hence so is every signed sum, hence so is every root.
\end{proof}

\subsection{Why complete splitting is not an algorithm}\label{sec:notalgo}

A reader meeting Theorem~\ref{thm:factorbase} for the first time may suspect it
trivializes relation search. It does not, and the reason is an asymmetry between
which variables are given and which are sought.

Theorem~\ref{thm:factorbase} is a statement about solving for the \emph{last}
coordinate given the first $m-1$. Its operational content is: choose
$P_1,\dots,P_{m-1} \in \mathcal{F}$; then the $2^{m-2}$ admissible values of the
remaining coordinate are $x(\pm P_1 \pm \cdots \pm P_{m-1})$, computable by
$O(2^{m-2})$ group operations. That is precisely the ``solve for the last
variable and test membership'' step that index calculus already performs, and it
is why the summation-polynomial system is usually posed in $m-1$ rather than $m$
unknowns.

Relation search asks the opposite question. The target $x(R)$ is \emph{fixed}, and
one must find $a_1,\dots,a_{m-1}$ \emph{in the factor base} with
$S_m(a_1,\dots,a_{m-1},x(R)) = 0$. Complete splitting in the last variable says
nothing about the existence or the cost of finding such a tuple: it does not make
the factor base membership test cheaper, it does not change how many tuples must
be tried, and it does not change the algebraic difficulty of the system when one
does pose it as a polynomial system. The cost of index calculus sits in that
search and in the subsequent linear algebra, neither of which
Theorem~\ref{thm:factorbase} touches.

What Theorem~\ref{thm:factorbase} does give is negative and methodological, and it
is the reason the theorem is stated at all:

\begin{corollary}[The census obstruction]\label{cor:census}
Any statistic of the form ``the factorization type, cycle type, or Frobenius
class of the fibre of $S_m$'' is \emph{identically constant} on the factor-base
locus, at every $m$ and for every curve. In particular, a Chebotarev-style census
of generic fibres --- sampling $(a_1,\dots,a_{m-1})$ uniformly from $\Fq^{m-1}$ ---
measures a quantity whose value on the locus where relation search operates is
known in advance and is the same for all curves.
\end{corollary}

This closes a natural but unproductive line of enquiry: one cannot find a curve
family with anomalous relation rates by looking for anomalous monodromy or
anomalous fibre statistics, because by Theorem~\ref{thm:main} there is no
anomalous monodromy to find, and by Theorem~\ref{thm:factorbase} the fibre
statistic is constant where it would have to matter.

%======================================================================
\section{Conservation of decomposition yield}\label{sec:conservation}

The second half of the synthesis is combinatorial and even more elementary. It is
recorded here because it constrains the same design space from the other side.

\begin{proposition}[Conservation]\label{prop:conservation}
Let $G$ be a finite abelian group of order $N$ and let
$\mathcal{D} \subseteq G \setminus \{0\}$ have $B$ distinct elements. For
$r \in G$ let $c_{\mathcal{D}}(r)$ be the number of size-$m$ multisets from
$\mathcal{D}$ summing to $r$. Then
\[
  \sum_{r \in G} c_{\mathcal{D}}(r) = \binom{B+m-1}{m},
  \qquad\text{hence}\qquad
  \mathbb{E}_{r}\bigl[ c_{\mathcal{D}}(r) \bigr] = \frac{1}{N}\binom{B+m-1}{m}
\]
for \emph{every} $\mathcal{D}$ of size $B$, independently of how $\mathcal{D}$ is
chosen.
\end{proposition}

\begin{proof}
Every size-$m$ multiset from $\mathcal{D}$ sums to exactly one element of $G$; count
the pairs (multiset, its sum) in two ways.
\end{proof}

\begin{corollary}\label{cor:conseq}
At matched factor-base size $B$:
\begin{enumerate}
\item The mean decomposition yield per target is a fixed function of $B$, $m$ and
  $N$. Any hypothesis asserting that some factor-base geometry has a mean-yield
  advantage --- or that such an advantage grows with $N$ --- is refuted by
  arithmetic, with no experiment.
\item Only the \emph{distribution} of $c_{\mathcal{D}}$ is free. Coverage, the
  fraction of targets admitting at least one decomposition, satisfies
  $\mathrm{cov} \le \min(1,\mu)$ with $\mu$ the mean, with equality iff no target
  has two decompositions. Additive structure that creates repeated sums therefore
  \emph{lowers} coverage at matched size.
\item Typing the decomposition (requiring the $i$-th summand to come from a
  sub-base $\mathcal{D}_i$ with $\sum_i B_i = B$) is a fixed penalty, not a lever:
  $\prod_i B_i \le (B/m)^m$, strictly below the untyped $\binom{B+m-1}{m} \approx
  B^m/m!$.
\end{enumerate}
\end{corollary}

Proposition~\ref{prop:conservation} is a one-line double count and no novelty is
claimed for it. It is included because, together with
\S\ref{sec:factorbase}, it delimits the design space from both sides, and because
it is easy to state hypotheses that it refutes without noticing.

\begin{remark}[What conservation does not say]
It does not say structured bases cannot beat random ones: the coverage headroom
up to $\min(1,\mu)$ is real and attainable (a Sidon/$B_h$ base attains it). It
says nothing about the \emph{cost of finding} a decomposition or about the linear
algebra --- which is where the cost of index calculus actually sits. And it is
entirely consistent with index calculus working: the Gaudry--Diem $1/n!$
decomposition probability over extension fields \emph{is} this conservation mean;
the extension-field advantage lives in the solve, not in the yield.
\end{remark}

%======================================================================
\section{Synthesis: which levers are fixed and which are free}\label{sec:synthesis}

\begin{center}
\begin{tabular}{@{}p{0.42\textwidth}p{0.52\textwidth}@{}}
\toprule
\textbf{Quantity} & \textbf{Status} \\
\midrule
Mean $m$-term decomposition yield at fixed factor-base size
  & \textbf{Fixed} by conservation (Prop.~\ref{prop:conservation}); identical
    for all factor bases of the same size. \\[2pt]
Galois group / monodromy of the summation cover
  & \textbf{Fixed} $= (\ZZ/2)^{m-2}$ for every curve, field and characteristic
    (Thm.~\ref{thm:main}); arithmetic $=$ geometric. \\[2pt]
Fibre factorization type on the factor-base locus
  & \textbf{Fixed} $=$ totally split, at every $m$
    (Thm.~\ref{thm:factorbase}). \\[2pt]
Fibre factorization type at a uniform random point
  & \textbf{Fixed} up to $O(m^2/q)$: totally split with density
    $2^{-(m-2)}$, else $2^{m-3}$ quadratics (Thm.~\ref{thm:frob},
    Cor.~\ref{cor:density}); trace cancels at first order. \\[2pt]
Square class of $\disc_T S_m$
  & \textbf{Fixed} $=$ square for $m \ge 4$; informative only at $m = 3$
    (Cor.~\ref{cor:disc}). \\[2pt]
Block systems of the fibre
  & \textbf{Fixed} $=$ subgroups of $(\ZZ/2)^{m-2}$; the recursion already
    exhibits the sign-subgroup family (Cor.~\ref{cor:blocks}). \\
\midrule
Coverage, and the multiplicity distribution
  & \textbf{Free}, bounded by $\min(1,\mu)$; attainable headroom is a constant
    factor, not an exponent. \\[2pt]
Relation rank and independence
  & \textbf{Free}; invisible to all of the above (the mean scores a
    rank-deficient base as tied with a random one). \\[2pt]
Recognizability: cost of testing factor-base membership
  & \textbf{Free}. \\[2pt]
Solve cost: the algebraic cost of the decomposition system
  & \textbf{Free}, and this is where the cost of index calculus sits. \\
\bottomrule
\end{tabular}
\end{center}

The headline is the separation. \emph{Generic-fibre arithmetic statistics and
factor-base relation statistics are different objects.} The first are governed by
Theorem~\ref{thm:main} and are the same for every curve; the second live on a
locus where, by Theorem~\ref{thm:factorbase}, the first are constant. A proposal
that hopes to raise the relation rate by changing curve or factor base must argue
in the ``free'' half of the table; the ``fixed'' half is closed by the results
above, and a proposal arguing there can be refuted on sight.

\begin{remark}[Scope of ``moves no exponent'']
The per-fibre statements above are at fixed $m$ and fixed field. A factor of $2$
applied once per summation fibre in a Gaudry/Diem-style decomposition over
$\FF_{q^n}$ with $n$ growing compounds as $2^{n-1}$ and \emph{would} move an
exponent. The claims of this section are made at fixed $m$ over a prime field and
must not be quoted without that scope.
\end{remark}

%======================================================================
\section{Computational verification}\label{sec:verify}

Every statement above that can be checked at finite scale was checked, on a
battery of $20$ curves over the primes $101, 103, 211, 307, 401, 1009$: twelve
of generic $j$; one $j = 0$ with $p \equiv 1 \bmod 3$ and one with
$p \equiv 2 \bmod 3$; one $j = 1728$ with $p \equiv 1 \bmod 4$ and one
supersingular with $p \equiv 3 \bmod 4$; and four with full rational $2$-torsion.
The implementation is self-contained (no computer-algebra system): summation
polynomials are built from \eqref{eq:s3} and \eqref{eq:recursion} by Sylvester
determinants recovered from $2^{m-2}+1$ evaluations, signed sums are computed by
elliptic-curve arithmetic over $\FF_{p^2}$, and factorization types are read off by
distinct-degree factorization. Seed $20260807$; single run, $68$ seconds.

\paragraph{The $m=3$ discriminant identity (Cor.~\ref{cor:disc}).}
Verified as an identity in $\ZZ[X_1,X_2,A,B]$ with exact integer coefficients ---
not as a congruence modulo a prime.

\paragraph{Root identification and the factorization law
(Thms.~\ref{thm:main}, \ref{thm:frob}).}
$18\,000$ uniform specializations at $m = 3,4,5$, of which $16\,737$ were good.
On all $16\,737$: the $2^{m-2}$ signed sums were distinct, none at infinity, and
each was a root of the computed polynomial ($16\,737/16\,737$); the factorization
type matched the prediction of Theorem~\ref{thm:frob} ($16\,737/16\,737$); no
irreducible factor of degree $>2$ ever appeared; and no fibre ever had a
\emph{mixed} factorization type. All $2^{m-2}$ Frobenius classes were realized at
each $m$, and each class produced exactly one factorization type ---
$1^{2^{m-2}}$ for the trivial class, $2^{2^{m-3}}$ for every other. The
discriminant square class was $+1$ on all $7324$ good samples at $m=4$ and all
$1627$ at $m = 5$, and mixed at $m = 3$ ($3828$ square, $3958$ non-square), as
Corollary~\ref{cor:disc} requires.

\paragraph{Factor-base locus (Thm.~\ref{thm:factorbase}).}
$18\,000$ specializations with every $a_i$ the $x$-coordinate of a rational
point, of which $16\,107$ were good. All $16\,107$ split completely, with all
$2^{m-2}$ roots in $\FF_p$ and equal to the signed sums. No counterexample at any
$m$ on any curve.

\paragraph{Exhaustive checks and the exact class count
(Cor.~\ref{cor:density}).}
All $p^{m-1}$ specializations were enumerated for $(m,p) \in \{3\}\times\{53,101,211\}$
and $\{4\}\times\{37,53\}$: $257\,061$ specializations, $224\,906$ of them good.
Zero factorization-type mismatches. For \emph{every} Frobenius class in every
cell, the identity
\[
  \underbrace{S^k N^{n-k} + N^k S^{n-k}}_{\text{closed form } \eqref{eq:classcount}}
  \;=\; \#\{\text{good } a \text{ in the class}\} +
        \#\{\text{degenerate } a \text{ in the class}\}
\]
held with residual exactly $0$. The degenerate locus matched
Proposition~\ref{prop:branch} exactly: for $m=3$ with $Z = 0$ it had exactly $p$
points ($101$ at $p=101$, $211$ at $p=211$ --- the diagonal $t_1 = t_2$), and for
$p = 53$ with $Z = 1$ exactly $157 = 53 + (2\cdot 53 - 1) - 1$ points.

\paragraph{Null controls.}
To show the instrument is discriminating rather than vacuous, the same pipeline
was run on two degree-$4$ families that are not summation polynomials.
NULL-A, a uniformly random degree-$4$ polynomial over $\FF_{401}$: of $599$
separable samples, only $13.7\%$ obeyed the dichotomy of
Theorem~\ref{thm:frob}, with types $1^13^1$, $4^1$ and $1^22^1$ all common.
NULL-B, the $S_4$ construction \eqref{eq:recursion} with the inner summation
polynomial replaced by a random quadratic --- the elliptic structure removed and
everything else held fixed: of $595$ separable samples, $48.9\%$ obeyed it, with
$4^1$ and $1^22^1$ both frequent. The factorization law was checked directly on
$16\,737 + 224\,906 = 241\,643$ good specializations of genuine summation
polynomials and held on every one; a law obeyed by $100\%$ of those and by
$13.7\%$ and $48.9\%$ of matched controls is not an artifact of the measurement.

%======================================================================
\section{Relation to prior work, and the limits of this novelty assessment}
\label{sec:novelty}

\paragraph{Classical, and claimed as such.}
Definition~\ref{def:sm} and the resultant recursion \eqref{eq:recursion} are
Semaev's \cite{semaev2004}. Lemma~\ref{lem:roots} --- the roots of
$S_m(a_1,\dots,a_{m-1},T)$ are $x$-coordinates of signed sums --- is an immediate
consequence of the defining property and is used routinely
\cite{gaudry2009,fghr2014,ky2015}; no novelty is claimed for it. The degree
$2^{m-2}$, the symmetry, and the absolute irreducibility of $S_m$ are standard
\cite{semaev2004,gaudry2009}, so Corollary~\ref{cor:irred} is a re-derivation, not
a new result. The operational fact behind \S\ref{sec:notalgo} --- that the last
coordinate is determined by group arithmetic --- is likewise standard practice.
Proposition~\ref{prop:conservation} is a double count.

\paragraph{What is offered as possibly new.}
Theorem~\ref{thm:main} (the Galois group for all $m$, with the identification of
the splitting field as $K(E^{m-1}/\Delta)$ and the coincidence of arithmetic and
geometric monodromy); Theorem~\ref{thm:frob} (the exact factorization dichotomy
over $\Fq$ and its dependence only on the class of the splitting-invariant
vector); Corollary~\ref{cor:disc} (the discriminant square class, isolating
$m=3$); and the synthesis of \S\ref{sec:synthesis} with
Corollary~\ref{cor:census}.

\paragraph{The limits of this assessment --- stated plainly.}
Prime-field Semaev index calculus has a substantial literature, and the
statements above are the kind that could plausibly be folklore among specialists
in summation polynomials or in function-field arithmetic. The following is the
exact extent of the prior-art search behind this paper, and it is not sufficient
to establish novelty:
\begin{itemize}
\item Targeted searches on the combination \{summation polynomial, Semaev\} with
  \{Galois group, monodromy, splitting field, factorization pattern, Kummer
  quotient, deck group, elementary abelian\} returned no work stating either
  Theorem~\ref{thm:main} or Theorem~\ref{thm:frob}.
\item \textbf{eprint.iacr.org was not reachable} from the environment in which
  this work was carried out (HTTP 403 to both the automated fetcher and a direct
  request). Since that is the primary venue for this literature --- including
  Semaev's original preprint --- a substantial fraction of the relevant corpus
  was \emph{not} inspected. arXiv was reachable and was inspected.
\item Consequently novelty is \textbf{not adjudicated} here in either direction.
  Before submission, a reader with access must check at minimum:
  \cite{semaev2004}, \cite{gaudry2009}, \cite{diem2011}, \cite{fghr2014},
  \cite{gg2014}, \cite{ky2015}, and the literature on point compression and
  Kummer-line arithmetic, where the quotient $E^n/\pm$ appears for unrelated
  reasons and where the deck-group observation may already be recorded.
\end{itemize}

\paragraph{Outstanding requirement before submission.}
Independent review of \S\ref{sec:main} by a reader experienced in function fields
or arithmetic geometry has \emph{not} been obtained. The proofs are short and the
computational evidence is strong, but the load-bearing steps ---
Lemma~\ref{lem:gamma}'s use of Artin's theorem for a quotient by a finite group
acting faithfully, and Lemma~\ref{lem:sep}'s properness argument in characteristic
$2$ where $E[2]$ is non-reduced --- deserve a second pair of eyes before this is
posted.

%======================================================================
\section{Open questions}

\begin{question}
The block systems of Corollary~\ref{cor:blocks} correspond to all subgroups of
$(\ZZ/2)^{m-2}$, while the resultant recursion exhibits only the family $H_I$. For
$m \ge 5$ the others are not visible in the recursion. Do they correspond to
resolvents with any computational meaning, or are they only group-theoretic
shadows?
\end{question}

\begin{question}
Theorem~\ref{thm:main} concerns the fibre in one variable with the others generic.
Relation search fixes the target and constrains all remaining variables to a
factor base. Is there a statistic on \emph{that} configuration --- the
intersection of $V(S_m)$ with a product of factor-base subschemes --- which is not
constant across curves? Corollary~\ref{cor:census} says the fibre statistics are
not; it does not say no statistic is.
\end{question}

\begin{question}
Corollary~\ref{cor:density} shows the trace cancels at first order in the split
density at every $m$. Is there a natural weighting of fibres under which the trace
survives?
\end{question}

%======================================================================
\begin{thebibliography}{9}

\bibitem{diem2011}
C. Diem.
\newblock On the discrete logarithm problem in elliptic curves.
\newblock \emph{Compositio Mathematica} 147(1):75--104, 2011.

\bibitem{fghr2014}
J.-C. Faugère, P. Gaudry, L. Huot, G. Renault.
\newblock Using symmetries in the index calculus for elliptic curves discrete
  logarithm.
\newblock \emph{Journal of Cryptology} 27(4):595--635, 2014.

\bibitem{fj2008}
M. D. Fried, M. Jarden.
\newblock \emph{Field Arithmetic}, 3rd ed.
\newblock Springer, 2008.

\bibitem{gg2014}
S. D. Galbraith, S. W. Gebregiyorgis.
\newblock Summation polynomial algorithms for elliptic curves in characteristic
  two.
\newblock In \emph{INDOCRYPT 2014}, LNCS 8885, pp. 409--427, 2014.

\bibitem{gaudry2009}
P. Gaudry.
\newblock Index calculus for abelian varieties of small dimension and the
  elliptic curve discrete logarithm problem.
\newblock \emph{Journal of Symbolic Computation} 44(12):1690--1702, 2009.

\bibitem{ky2015}
M. Kosters, S. L. Yeo.
\newblock Notes on summation polynomials.
\newblock arXiv:1503.08001, 2015.

\bibitem{semaev2004}
I. Semaev.
\newblock Summation polynomials and the discrete logarithm problem on elliptic
  curves.
\newblock IACR Cryptology ePrint Archive, Report 2004/031, 2004.

\bibitem{silverman}
J. H. Silverman.
\newblock \emph{The Arithmetic of Elliptic Curves}, 2nd ed.
\newblock GTM 106, Springer, 2009.

\end{thebibliography}

\end{document}
